\documentclass[draft,a4paper,14pt]{extarticle} % 14й шрифт
%%% Преамбула %%%

\usepackage[cm-default]{fontspec} % XeTeX
\usepackage{xunicode} % Unicode для XeTeX
\usepackage{xltxtra}  % Верхние и нижние индексы
\usepackage{pdfpages} % Вставка PDF
\usepackage{bm}

\usepackage{listings} % Оформление исходного кода
\lstset{
    basicstyle=\small\ttfamily, % Размер и тип шрифта
    breaklines=true, % Перенос строк
    tabsize=2, % Размер табуляции
    literate={--}{{-{}-}}2 % Корректно отображать двойной дефис
}
% Русский язык
\usepackage{polyglossia}
\setdefaultlanguage{russian}
\setotherlanguage{english}

% Шрифты, xelatex
\defaultfontfeatures{Ligatures=TeX}
\setmainfont[Ligatures=TeX]{Times New Roman}
\newfontfamily\cyrillicfont{Times New Roman}[Script=Cyrillic]
% \setsansfont{Liberation Sans} % Альтернатива нью роману, если его нет
% \setmonofont{FreeMono} % Моноширинный шрифт для оформления кода
\newfontfamily{\cyrillicfonttt}{FreeMono}


\usepackage{amssymb,amsfonts,amsmath} % Математика
\numberwithin{equation}{section} % Формула вида секция.номер

% \let\oldequation=\equation
% \let\endoldequation=\endequation
% \renewenvironment{equation}{\begin{oldequation}}{\end{oldequation}}
\newenvironment{eqseries}
{\vskip 0pt}
{\vskip 24pt}

\usepackage{enumerate} % Тонкая настройка списков
\usepackage{indentfirst} % Красная строка после заголовка
\usepackage{float} % Расширенное управление плавающими объектами
\usepackage{multirow} % Сложные таблицы
\usepackage{adjustbox}
\usepackage{lscape}
\usepackage{rotating}

\usepackage{pdfcomment}     % тудущки и комменты
\pdfcommentsetup{final} % set to 'final' to prevent all comments from showing up   draft
\pdfcommentsetup{color={1.0 1.0 0.0}, open=true}

% Пути к каталогам с изображениями
\usepackage{graphicx} % Вставка картинок и дополнений
\graphicspath{{pics/}{images/userguide/}{images/testing/}{images/infrastructure/}{extra/}{extra/drafts/}}

% Формат подрисуночных записей
\usepackage{chngcntr} 
\counterwithin{figure}{section}

% Гиперссылки
\usepackage{hyperref}
\hypersetup{
    colorlinks, urlcolor={black}, % Все ссылки черного цвета, кликабельные
    linkcolor={black}, citecolor={black}, filecolor={black},
    pdfauthor={Author},
    pdftitle={Title}
}
% Оформление библиографии и подрисуночных записей через точку
\makeatletter
\renewcommand*{\@biblabel}[1]{\hfill#1.}
\renewcommand*\l@section{\@dottedtocline{1}{1em}{1em}}
\renewcommand{\thefigure}{\thesection.\arabic{figure}} % Формат рисунка секция.номер
\renewcommand{\thetable}{\thesection.\arabic{table}} % Формат таблицы секция.номер
\def\redeflsection{\def\l@section{\@dottedtocline{1}{0em}{10em}}}
\makeatother

\renewcommand{\baselinestretch}{1.4} % Полуторный межстрочный интервал
\parindent 1.27cm % Абзацный отступ

\sloppy             % Избавляемся от переполнений
\hyphenpenalty=1000 % Частота переносов
\clubpenalty=10000  % Запрещаем разрыв страницы после первой строки абзаца
\widowpenalty=10000 % Запрещаем разрыв страницы после последней строки абзаца

% Отступы у страниц
\usepackage{geometry}
\geometry{left=3cm}
\geometry{right=1.5cm}
\geometry{top=2cm}
\geometry{bottom=2cm}

% Списки
\usepackage{enumitem}
\setlist[enumerate,itemize]{leftmargin=0mm,itemindent = 20mm} % Отступы в списках

\makeatletter
    \AddEnumerateCounter{\asbuk}{\@asbuk}{м)}
\makeatother
\setlist{nolistsep} % Нет отступов между пунктами списка
\renewcommand{\labelitemi}{--} % Маркет списка --
\renewcommand{\labelenumi}{\asbuk{enumi})} % Список второго уровня
\renewcommand{\labelenumii}{\arabic{enumii})} % Список третьего уровня
                    % Глубина оглавления, до subsubsection
\usepackage{scrextend}


% Нумерация страниц по центру снизу
\usepackage{fancyhdr}
\pagestyle{plain}
\fancyhf{}
\fancyhead[R]{\textrm{\thepage}}
\fancyheadoffset{0mm}
\fancyfootoffset{0mm}
\setlength{\headheight}{17pt}
\renewcommand{\headrulewidth}{0pt}
\renewcommand{\footrulewidth}{0pt}
\fancypagestyle{plain}{ 
    \fancyhf{}
    \rhead{\thepage}
}
% Содержание
\usepackage{tocloft}
\renewcommand{\cftaftertoctitle}{
    \thispagestyle{plain}}     % Убирает номер страницы
\renewcommand{\cfttoctitlefont}{\hspace{0.38\textwidth}\MakeTextUppercase} % СОДЕРЖАНИЕ
\renewcommand{\cftsecfont}{\hspace{0pt}}            % Имена секций в содержании не жирным шрифтом
\renewcommand\cftsecleader{\cftdotfill{\cftdotsep}} % Точки для секций в содержании
\renewcommand\cftsecpagefont{\mdseries}             % Номера страниц не жирные
% \renewcommand{\cftparskip}{0pt}
\setlength{\cftbeforesecskip}{0pt}
\setlength{\cftsecindent}{5mm}
\setlength{\cftsubsecindent}{10mm}
\setcounter{tocdepth}{3}        

% Список иллюстративного материала
\renewcommand{\cftloftitlefont}{\hspace{0.17\textwidth}\MakeTextUppercase}
\renewcommand{\cftfigfont}{Рисунок }
\addto\captionsrussian{\renewcommand\listfigurename{Список иллюстративного материала}}
% Список табличного материала
\renewcommand{\cftlottitlefont}{\hspace{0.2\textwidth}\MakeTextUppercase}
\renewcommand{\cfttabfont}{Таблица }
\addto\captionsrussian{\renewcommand\listtablename{Список табличного материала}}


% GOST - biblatex
% \usepackage[% 
% bibstyle=gost-numeric, %подключение одного из четырех главных стилей biblatex-gost 
% sorting=none, %тип сортировки в библиографии
% doi=false,  
% ]{biblatex}
% \usepackage[style=gost2003.bst , backend=bibtex8]{biblatex}
\bibliographystyle{gost2003}
\renewcommand\refname{книги}
% \addbibresource{refs.bib}
% \bibliography{refs.bib}

% Формат подрисуночных надписей
\RequirePackage{caption}
\DeclareCaptionLabelSeparator{defffis}{ -- } % Разделитель
\captionsetup[figure]{justification=centering, labelsep=defffis, format=plain, belowskip=-5pt, aboveskip=20pt, font={singlespacing}} % Подпись рисунка по центру
\captionsetup[table]{justification=raggedright, labelsep=defffis, format=plain, belowskip=0pt, aboveskip=20pt, singlelinecheck=false} % Подпись таблицы слева
\addto\captionsrussian{\renewcommand{\figurename}{Рисунок}} % Имя фигуры

% Пользовательские функции
\newcommand{\addimg}[4]{ % Добавление одного рисунка
    \begin{figure}
        \centering
        \includegraphics[width=#2\linewidth]{#1}
        \caption{#3} \label{#4}
    \end{figure}
}
\newcommand{\addimghere}[4]{ % Добавить рисунок непосредственно в это место
    \begin{figure}[H]
        \centering
        \includegraphics[width=#2\linewidth]{#1}
        % \medskip
        \caption{#3} \label{#4}
        % \medskip
    \end{figure}
}
\newcommand{\addtwoimghere}[6]{ % Вставка двух рисунков
    \begin{figure}[H]
        \centering
        \includegraphics[width=#2\linewidth]{#1}
        \hspace{1cm}
        \includegraphics[width=#4\linewidth]{#3}
        \caption{#5} \label{#6}
    \end{figure}
}
\newcommand{\addimgapp}[2]{ % Это костыль для приложения Б
    \begin{figure}[H]
        \centering
        \includegraphics[width=1\linewidth]{#1}
        \caption*{#2}
    \end{figure}
}

% Заголовки секций в оглавлении в верхнем регистре
\usepackage{textcase}
\makeatletter
\let\oldcontentsline\contentsline
\def\contentsline#1#2{
    \expandafter\ifx\csname l@#1\endcsname\l@section
        \expandafter\@firstoftwo
    \else
        \expandafter\@secondoftwo
    \fi
    {\oldcontentsline{#1}{#2}}          % {\oldcontentsline{#1}{\MakeTextUppercase{#2}}} 
    {\oldcontentsline{#1}{#2}}
}
\makeatother

% Оформление заголовков
\usepackage[compact,explicit]{titlesec}
\titleformat{\section}{\large\bfseries}{}{0mm}{\centering{\thesection\quad\textbf{#1}} }  % \titleformat{\section}{}{}{12.5mm}{\centering{\thesection\quad\MakeTextUppercase{#1}}\vspace{1.5em}}
\titleformat{\subsection}[block]{\bfseries}{}{12.5mm}{\thesubsection\quad\textbf{#1}}          %
\titleformat{\subsubsection}[block]{\vspace{1em}\normalsize}{}{12.5mm}{\thesubsubsection\quad#1\vspace{1em}}
\titleformat{\paragraph}[block]{\normalsize}{}{12.5mm}{\MakeTextUppercase{#1}}


% Секции без номеров (введение, заключение...), вместо section*{}
\newcommand{\anonsection}[1]{
    \phantomsection % Корректный переход по ссылкам в содержании
    \paragraph{\centerline{{#1}}\vspace{0em}}
    \addcontentsline{toc}{section}{\hspace{-5mm}\uppercase{#1}}
}

% Секция для аннотации (она не включается в содержание)
\newcommand{\annotation}[1]{
    \paragraph{\centerline{{#1}}\vspace{1em}}
}

% Секция для списка иллюстративного материала
\newcommand{\lof}{
    \phantomsection
    \listoffigures
    \addcontentsline{toc}{section}{\listfigurename}
}

% Секция для списка табличного материала
\newcommand{\lot}{
    \phantomsection
    \listoftables
    \addcontentsline{toc}{section}{\listtablename}
}

% Секции для приложений
\newcommand{\appsection}[1]{
    \phantomsection
    \paragraph{\centerline{{#1}}}
    \addcontentsline{toc}{section}{\uppercase{#1}}
}

% Библиография: отступы и межстрочный интервал
% \makeatletter
% \defbibenvironment{thebibliography}
% {
%     \list{\@biblabel{\@arabic\c@enumiv}}
%         {\settowidth\labelwidth{\@biblabel{#1}}
%         \leftmargin\labelsep
%         \itemindent 16.7mm
%         \@openbib@code
%         \usecounter{enumiv}
%         \let\p@enumiv\@empty
%         \renewcommand\theenumiv{\@arabic\c@enumiv}}
%         }
%     \setlength{\itemsep}{0pt}
%     }
% \makeatother

\setcounter{page}{4} % Начало нумерации страниц % Подключаем преамбулу

%%% Начало документа
\begin{document}

\includepdf{include/titlepage}
%\includepdf{pz} % Пояснительная записка
\includepdf[pages=1-2]{include/task} % Задание на диплом 
\includepdf[pages=1-2]{include/referat}

\setcounter{page}{0}
\hyphenchar\font=-1 % disable hyphen
\tableofcontents % Содержание 
\hyphenchar\font=`\- % reset hyphen
\thispagestyle{empty}
\clearpage
\setcounter{page}{6} % Начало нумерации страниц

\anonsection{введение}
За последние годы был совершён существенный прогресс в доступности и точности сенсоров, позволяющих мобильным роботам 
осуществлять оценку окружающего пространства. Такие информационно-измерительные устройства как лидары, сонары и стереокамеры
 стали основной опорой для алгоритмов автономной навигации и локализации. Тем не менее в роботах по-прежнему 
присутствуют телевизионные системы, так как они дают наиболее легко воспринимаемую информацию для оператора в случаях, когда 
его вмешательство необходимо. 
Число и расположение видеокамер может быть различным. Бортовые телевизионные системы классифицируются по различным 
критериям\cite{varlashin}:
\begin{itemize} 
    \item по количеству камер (монокулярные или многокамерные).
    \item по функциональному назначению:
    \begin{enumerate}[leftmargin=12mm]  %\setlength{\leftmargin}{10mm}
        \item одиночные камеры;
        \item PTZ-камеры (англ. Pan-Tilt-Zoom);
        \item стереокамеры;
        \item системы кругового обзора;
        \item комбинированные системы.
    \end{enumerate}
\end{itemize}

Для наиболее эффективного покрытия телевизионной системой максимального объёма окружающего пространства в ней могут использоваться 
камеры "рыбий глаз" (англ. fisheye-camera), позволяющие одним кадром покрыть угловое поле свыше $180^\circ$.  
Набор таких камер может составлять систему кругового обзора, позволяющую оператору видеть не только в любом направлении,                
но даже с видом от третьего лица \cite{birdeye}. 
    
В случае автономных мобильных роботов телевизионный круговой обзор применяется лишь по необходимости, но при этом может быть весьма 
дорогостоящими и занимать место в корпусе. При этом в случае наличия пересечений полей зрения камер, входящих в систему, теоретически возможно 
применение принципа стереозрения для построения карты окружающего пространства. Однако значительные радиальные искажения изображения, вызванные 
особенностями используемых объективов, не позволяют использовать известные алгоритмы стереозрения без предварительной обработки изображений.                                   
 
  Реализация системы стереозрения на основе ортогонально ориентированных камер, способной эффективно задействовать особенности 
  сверхширокоугольных объективов, позволит дать робототехнически системам новый способ получать информацию об окружении.

Целью работы является разработка и исследование точности системы стереозрения, основанной на ортогонально расположенных сверхширокоугольных камерах.

В ходе работы решаются следующие задачи:
\begin{itemize}     % TODO: уточнить
    \item обзор современных систем стереозрения, использующих изображения с широкоугольных камер,
    \item обзор моделей искажений, применяющихся в моделировании камер с объективами "рыбий глаз",
    \item разработка алгоритма устранения искажений fisheye-изображения в  области кадра,                                 
    \item теоретическое определение точности калибровочных моделей камер,
    \item практическое определение точности системы.  
\end{itemize}


\clearpage
\section{Аналитический обзор систем стереозрения}

\subsection{Стереозрение}
\label{stereovision}
Система стереозрения состоит из двух камер, наблюдающих сцену с разных точек. 
Фундаментальная основа принципа заключается в предположении, что каждой точке в пространстве соответствует уникальная пара пикселей на снимках с двух камер.  

При этом к камерам предъявляются некоторые требования \cite{rusoverview}:   % не уверен, что это надо цитировать
\begin{itemize}
	\item Камеры откалиброваны. Это значит, что известны внутренние (оптические) и внешние (расположение камер в пространстве) параметры камер. 
	\item Ректификация. Подразумевает выравнивание изображения с обеих камер по строкам.  % Мб подробнее расписать  
	\item Ламбертовость поверхностей. Означает независимость освещения наблюдаемых поверхностей от угла зрения. 
\end{itemize}

Таким образом, соблюдение указанных выше требований позволяет использовать следующий геометрический принцип. При наличии двух камеры, как изображено 
на рисунке \ref{pic:epipol} \cite{Hartley2004}, где $C$ — центр первой камеры, $C'$ — центр второй камеры, точка пространства $X$  
проецируется в точки $x$ на плоскости изображения левой камеры и $x'$ на плоскости изображения правой камеры. Прообразом точки $x$ на изображении левой 
камеры является луч $xX$. Этот луч проецируется на плоскость второй камеры в прямую $l'$, называемую эпиполярной линией. Образ точки $X$ на плоскости 
изображения второй камеры обязательно лежит на этой линии.


В результате каждой точке $x$ на изображении левой камеры соответствует эпиполярная линия $l'$ на изображении правой камеры. При этом соответствие для $x$ на 
изображении правой камеры может лежать только на соответствующей эпиполярной линии. Аналогично, каждой точке $x'$ на правом изображении соответствует 
эпиполярная линия $l$ на левом. Поиск соответствий между точками $x$ и $x'$ называют стереосопставлением.

\addimghere{epipolar geometry}{0.5}{Эпиполярная геометрия}{pic:epipol}

Далее с помощью точек $x$ и $x'$ возможно посчитать смещения каждого пикселя одного изображения относительно другого, что позволяет построить карту смещений. 
Очевидно, что смещения будут подсчитаны только для точек, видимых обеими камерами. Карта смещений же приводится далее либо к облаку точек, либо к карте глубины. 
Стереосистемы, реализующие этот принцип, называют пассивными. Они являются самыми простыми и часто используются, так как для их создания достаточно лишь
двух зафиксированных камер. Однако пассивные системы опираются целиком на видимый свет, что ограничивает их применимость в экстремальных условиях освещённости. %\cite{find:passive_perfomance}. 

Одну из камер можно заменить источником  света, освещающим одну или несколько точек поверхности световым лучом или специальным структурированным шаблоном освещения.  % TODO: цитирование страницы  529
Так как структура шаблона и направление его  лучей заранее известны, камера может проводить оценку формы объектов в кадре по искажениям шаблона\cite{shapiro}. 
Структурная подсветка тоже является распространённым методом, реализованным во многих коммерческих сенсорах, благодаря совмещению высокой производительности 
и  низкой цены \cite{struct_light}.  

Совмещение пассивного стереозрения и структурированной подсветки позволяет улучшить качество стереосопоставления, особенно для поверхностей со слабо выраженной текстурой
и при низкой освещённости. Однако инфракрасный узор подсветки весьма ограничен в дальности и не виден при сильном солнечном освещении, что затрудняет использование этих    % FIXME: криво ппц
сенсоров на открытых пространствах. Системы с подобной технологий называют активными \cite{kinect_perf}. 

\addimghere{stereo_types}{1.0}{Виды организации стереосистем: а)~пассивная стереосистема; б)~активная стереосистема; в)~подсветка структирированным светом }{pic:stereo_methods}
%\pdfmargincomment{просится какой-то вывод}

Описанные подходы хорошо работают для обычных камер с незначительными искажениями, где легко выполнить ректификацию. 		%\pdfmargincomment{и сюда сразу модель обскуры или потом оставить?}
В случае  же камер с широкоугольными (от 60° до 100°), особоширокоугольными (от 100° до 180°) и сверхширокоугольными (свыше 180°)  \cite{camera_class} объективами
 в изображение вносятся существенные искажения, которые затрудняют поиск соответствий. На 
рисунке \ref{pic:fy_epipol} вручную сопоставлены одни и те же точки на фрагментах двух изображений со сверхширокоугольных камер,
смотрящих в одном направлении. Видно, что пары признаков больше не лежат на одной горизонтальной прямой.  
% Учитывая распространённость таких  камер в робототехнике, задача реализации систем стереозрения на их основе является актуальной.    

\addimghere{not_epipolar}{0.7}{Соответствия на снимках с fisheye-камер}{pic:fy_epipol}   % TODO: без поворота

Существует несколько основных направлений исследований в области методов стереозрения, применимых к таким камерам:
\begin{itemize}
\item исследования стереопар с двумя камерами со сверхширокоугольными объективами, направленными в одну сторону и имеющими параллельные оптические оси \cite{parallel_fy, sweep_net};
\item исследования стереопар с двумя камерами со сверхширокоугольными объективами, направленными в противоположные стороны и имеющими коллинеарные оптические оси;
\item исследования стереопар с двумя камерами со сверхширокоугольными объективами, имеющие ортогонально ориентированными оптические оси, лежащие в одной плоскости \cite{ortho_fy};
\end{itemize}

Первый подход является одним из самых распространённых. Последние работы в этой области предлагают \cite{direct_neuro_stereo} отказаться от типичных 
для стереозрения этапов устранения искажений и ректификации и извлекать информацию о глубине по двум снимкам сверхширокоугольных камер без 
предварительной обработки. Для производства карт глубины используется свёрточная нейронная сеть, что требует существенных вычислительных 
мощностей - для достижения производительности в реальном времени разработчикам понадобилось использовать настольный компьютер c мощным графическим ускорителем. 
В мобильном автономном роботе аналогичный по производительности вычислитель разместить может быть проблематично. 

С помощью второго подхода может быть реализована, например, кольцевая область стереозрения с вертикальным  %Метод с разворотом камер на 180  ...
полем зрения $65^\circ$ \cite{omni_stereo}. Для этого две $245^\circ$ камеры закреплены на противоположных концах жёсткого стержня 
и направлены друг на друга. Это позволяет достигнуть панорамного обзора глубины с качеством, достаточным для осуществления автономной навигации и
локализации БПЛА\cite{omni_copter}, но конструктивно такая схема расположения камер имеет смысл только для летательных аппаратов.  

Третий подход пока самый малоисследованный. В рамках него Чжан (Zhang) и другие разработали особую систему стереозрения (the special stereo vision 
system), использующую модуль из четырёх камер с полем зрения $185^\circ$. Размещены камеры в одной плоскости под углом в $90^\circ$ друг к другу, 
как изображено на рисунке \ref{pic:ssvs} \cite{zhang_system}. Авторы отразили в работе калибровку разработанной ими системы и устранение искажений, 
но не  продемонстрировали результаты оценки глубины и не проанализировали точность метода. 

\addimghere{ssvs}{0.6}{Модуль камер}{pic:ssvs}

\clearpage
\vspace{\baselineskip}
            % Обзор систем
\subsection{Примеры применения сверхширокоугольных камер в робототехнике }

Рассмотрим применение и расположение сверхширокоугольных камер в некоторых существующих роботах и 
технических системах.

Камеры со сверхшироким полем зрения уже давно используются в конструкции марсоходов NASA. Типичная система 
камер ровера помимо прочего состоит из пары NavCam и трёх пар HazCam \cite{opportunity}. NavCam - это камеры, использующиеся для навигации,
 панорамных снимков и определения препятствий. Они размещены высоко, обладают узким диагональным полем зрения примерно $67^\circ$ (в марсоходе 
Perseverance повышено до $120^\circ$ \cite{perseverance}) и большим фокусным расстоянием. HazCam - это камеры, направленные на грунт перед марсоходом для обнаружения 
 препятствий в ближнем поле. Они оснащены fisheye-линзами с полем зрения примерно $120^\circ * 120^\circ$ (в марсоходе Perseverance изменены на $136^\circ * 102^\circ$).
Расположение камер изображено на рисунке \ref{pic:perseverance}. Каждая пара камер используется для построения карты глубины в соответствующей ей дальности 
и области пространства вокруг ровера. 

\addimghere{pics/perseverence}{0.7}{Расположение камер на марсоходе Perseverance; Right NavCam, Left NavCam - соответсвтенно правая и левая камера NavCam; RA и RB - правая пара камер HazCam; LA и LB - левая пара камер HazCam \ref{pic:perseverance}}{pic:perseverance}

Робот-доставщик "Ровер R3"  компании Яндекс имеет на борту 5 сверхширокоугольных камер, 
размещённые спереди, сзади и по бортам корпуса \cite{yandex_rover}. Разработчики не раскрывают их назначения, но наиболее 
вероятные варианты - визуальная одометрия и обнаружение и слежение за объектами. % FIXME: можно ли так? 
На рисунке \ref{pic:4cam_system}.а показано реальное положение сенсоров робота, а на рисунке \ref{pic:4cam_system}.б 
изображена схема зон перекрытия сверхширокоугольных камер. 

\addimghere{rover_with_zones}{0.8}{а)~положение сенсоров на роботе; б)~зоны перекрытия полей зрения}{pic:4cam_system}
 
Кроме того многие автопроизводители в своих автомобилях реализуют системы кругового обзора для помощи в парковке или 
системы автономного вождения с применением fisheye-камер. Их положение у разных производителей и моделей может отличаться, 
но общая тенденция - размещать камеры, обеспечивая максимальное покрытие окружающего пространства. Пример положения и 
углов зрения камер автомобиля Tesla представлен на рисунке \ref{pic:tesla_cams}. 

\addimghere{tesla_cameras}{0.8}{Камеры автомобиля Tesla}{pic:tesla_cams}

Рассмотренные примеры показывают системы, уже обладающие сверхширокоугольными камерами. Внедрение в них системы стереозрения 
позволило бы получить дополнительный источник информации о окружающем пространстве и/или уменьшить количество используемых камер. 
Располагаются камеры под углом, близким к $90^\circ$, что не позволяет использовать существующие методы стереозрения и 
мотивирует к разработке системы стереозрения именно для такого случая.       
%\vspace{\baselineskip}
\subsection{Выводы по первому разделу}

% Обзор современных моделей искажений сверхширокоугольных камер позволил выбрать 
%  ряд наиболее точных и удобных для калибровки для дальнейшего сравнений. 
Рассмотрен принцип и классификация способов стереозрения. 
Осуществлён обзор существующих систем стереозрения, применяющих 
сверхширокоугольные камеры с целью ознакомления с мировым опытом. 
Рассмотрены основные тенденции в расположении камер и выборе линз 
в существующих мобильных роботах.

Принято решение разрабатывать и тестировать системы стереозрения,
основанную на сверхширокоугольных камерах, расположенных ортогонально.  

% Принято решение разрабатывать и тестировать 
% предлагаемую систему стереозрения с применением виртуального  моделирования. 

% Для моделирования выбрана среда разработки Unity. Для обработки изображений
% с камер выбрана библиотека OpenCV для языка программирования C++.              % Выводы по разделу 

\newpage
\section{Система стереозрения}
\subsection{Описание системы стереозрения}

Предлагаемая система стереозрения состоит из двух камер с объективами типа "рыбий глаз" $\geqslant180^\circ$,
расположенных ортогонально так, что  две камеры имеют область пересечения полей зрения. В пространстве, наблюдаемом 
сразу  двумя камерами проводится триангуляция и получение информации об объёме после этапа устранения искажений.  % FIXME: Всё описание просто бггг

% Рассмотрим организацию системы на примере робота-доставщика "Ровер R3"  компании Яндекс, который имеет на борту 4 сверхширокоугольные камеры, 
% размещённые спереди, сзади и по бортам корпуса \cite{yandex_rover}, что соответствует описанию системы. 
% На рисунке \ref{pic:4cam_system}, а показано реальное положение камер  робота и их зон перекрытия, в которых может 
%  обеспечивается получение информации о глубине при использовании описываемой системы. Рисунок \ref{pic:4cam_system}, б
% демонстрирует эквивалентную по горизонтальному  покрытию схема при использовании обычных камер 
%  с углом обзора $90^\circ$. 
 
% %\addtwoimghere{Group 1}{0.4}{Group 2}{0.4}{Сравнение систем стереозрения}{pic:4cam_system}  % TODO: разобраться с масштабом
% \addimghere{group12}{0.7}{Геометрическая модель бинокулярной системы стереозрения}{pic:4cam_system}

% Как можно заметить, системе на основе обычных камер нужно в два раза больше сенсоров, чтобы достичь той же зоны покрытия 
% по горизонтали.  Кроме того, традиционная система имеет меньшую зону покрытия по вертикали и не обеспечивает полный 
% панорамный обзор. Все эти факторы делают систему стереозрения на основе ортогонально расположенных сверхширокоугольных камер 
% более предпочтительной для применения в робототехнике.  % FIXME: уточнить. не во всей не всегда. сскорее выгодной, но  слово не очень

Применение существующих алгоритмов стереосопоставления предполагает наличие стереопары, удовлетворяющей условиям, описанным в секции \ref{stereovision}.
Соответственно, необходимо 
По двум ортогональным fisheye-снимкам такую стереопару можно получить, введя в систему  для каждой сверхширокоугольной камеры виртуальную(ые) камеру-обскуру и направив 
её в сторону пересечения полей зрения, как если бы это была пара обычных камер. Процесс формирования виртуальной камеры-обскуры и 
алгоритм устранения искажений более подробно описаны в секции \ref{dewarping}.

Далее для упрощения рассмотрения системы будет считаться, что оптические оси всех камер находятся в одной плоскости, 
а на ориентацию виртуальных камер влияет только угол $ \beta $ поворота в этой плоскости. 

На рисунке \ref{pic:2cam_scheme} изображён простейший вариант системы с двумя камерами под углом $90^\circ$, который можно  считать частью
схемы, представленной на рисунке \ref{pic:4cam_system}, или модуля, изображённого  на рисунке \ref{pic:ssvs}. 

\addimghere{sample_simple2cam}{0.7}{Геометрическая модель бинокулярной системы стереозрения}{pic:2cam_scheme} %TODO: перерисовать схему?

  Здесь область пересечения полей зрения камер $C_0$ и $C_1$ обозначенная красным.
Эта область эквивалентна области пересечения полей зрения двух камер с полями зрения $90^\circ$ (обозначены
оранжевым), повёрнутых на $\beta_0 = \beta_1  = 45^\circ$ в сторону области интереса. Тогда $B$ - база стереопары.
Примеры изображений, полученных в такой конфигурации, приведён на рисунке \ref{pic:dewarped_exmples}. 

\addimghere{4pic_example}{0.7}{Пример исходных изображений с отмеченной областью интереса и снимков виртуальной стереопары}{pic:dewarped_exmples}


      % Описание системы: Какие этапы, элементы и тд
\subsection{Алгоритм устранения искажений}
\label{dewarping}

Распространённые библиотеки машинного зрения предлагают готовые к использованию классы и функции, 
позволяющие устранять искажения fisheye-объективов. Например, в составе MATLAB Computer Vision Toolbox присутствует 
инструмент Camera Calibrator, позволяющий в несколько простых этапов  произвести оценку внутренних и внешних параметров
 камеры. Пользователю надо лишь  загрузить набор изображений, содержащих в себе калибровочный  узор, выбрать модель камеры 
(стандартная или сверхширокоугольная) и запустить автоматический процесс калибровки. По завершению процесса программа 
показывает точность калибровки и может отображать скорректированные изображения. Результаты можно экспортировать в рабочую 
область MATLAB для дальнейшей обработки изображений. Интерфейс программы представлен на рисунке \ref{pic:calibrator}. 

 \addimghere{calibrator}{0.8}{Интерфейс приложения MATLAB Camera Calibrator}{pic:calibrator} 

Однако на практике этот и другие подобные методы весьма ограничены. И MATLAB, и OpenCV устраняют искажения  строго 
в центральной части изображения, не позволяя выбирать  желаемое направление обзора.  Это ограничивает применимость 
существующих инструментов только для копланарного расположения  камер. Поэтому для реализации предлагаемой системы 
стереозрения  был разработан собственный алгоритм устранения искажений. Схема геометрического принципа, лежащего в
 основе этого алгоритма, представлена на рисунке \ref{pic:sweeping}.

\addimghere{projection_sweeping}{0.65}{Схема принципа устранения искажений}{pic:sweeping} 

Цель алгоритма - найти, куда на выходном изображении  с  перспективной проекцией проектируются все пиксели из 
выбранного участка входного сверхширокоугольного изображения. Сделать это можно, выполнив сначала обратное 
преобразование для каждого пикселя fisheye-снимка по модели сверхширокоугольной камеры, а затем прямое преобразование 
по модели камеры-обскуры. Однако такой подход приведёт к возникновению дефектов из-за несовпадения % TODO:  скорее всего наврал. Надо проконсультироваться 
частот дискретизации двух изображений.  Поэтому разработанный алгоритм сначала выполняет обратное преобразование  для 
каждого пикселя итогового изображения $\nu$, находя таким образом соответствующую ему точку в системе координат 
камеры $({X_c, Y_c, Z_c})$:
\vskip 12pt
\begin{equation}
    \label{eq:uv_to_xyz}
    \left[\begin{matrix}x_c\\y_c\\z_c\\\end{matrix}\right] = \left[\begin{matrix} {u*z_c}/f \\  {v*z_c}/f \\ z_c \\\end{matrix}\right],
\end{equation} 
\vskip 24pt

где  $f$ - фокусное расстояние. 

Набор таких точек формирует прямоугольную область $\nu_p$ с центром в точке $O_p$ и является плоскостью изображения
 виртуальной камеры-обскуры  с оптической осью $Z_p$. Поворот точек, входящих в $\nu_p$, с помощью матрицы вращения 
 $\bm{R}$ образует $\nu'_p$ и позволяет таким образом задать направление обзора и ориентацию виртуальной камеры. 
\begin{eqseries}
    \begin{equation}
        \label{eq:R}
        \bm{R} = \left[\begin{matrix}\cos{\alpha}&-\sin{\alpha}&0\\\sin{\alpha}&\cos{\alpha}&0\\0&0&1\\\end{matrix}\right]\left[\begin{matrix}1&0&0\\0&\cos{\beta}&\sin{\beta}\\0&-\sin{\beta}&\cos{\beta}\\\end{matrix}\right]\left[\begin{matrix}\cos{\gamma}&0&-\sin{\gamma}\\0&1&0\\\sin{\gamma}&0&\cos{\gamma}\\\end{matrix}\right],
    \end{equation} 
    \begin{equation}
        \label{eq:sweeped}
        \left[\begin{matrix}x'_p\\y'_p\\z'_p\\\end{matrix}\right] = \left[\begin{matrix}x_p\\y_p\\z_p\\\end{matrix}\right] \bm{R}.
    \end{equation}  
\end{eqseries}

где координаты с индексом $p$ -  координаты точек соответствующей плоскости $\nu$ в системе камеры, 

\qquad $\alpha, \beta, \gamma$ - углы Эйлера, управляющие  ориентацией виртуальной камеры-обскуры. 

Тогда обратная fisheye-проекция точек из $\nu'_p$ позволяют получить область $\nu_i$ исходного изображения. 
Таким образом, зная связь всех точек из $\nu$ и $\nu_i$, возможно переносить информацию о цвете из 
исходного изображения на итоговое.  

Данный принцип устранения искажений может быть применён к любой модели сверхширокоугольной камеры, если
она обладает явной обратной проекцией. Очевидно, весь процесс преобразования $\nu \rightarrow \nu_i$ требует 
выполнения существенного количества математических операций, что негативно сказывается на скорости, с 
которой алгоритм может обрабатывать изображения в реальном времени. Однако при неизменных параметрах 
модели алгоритм достаточно выполнить лишь один раз, записав результат в таблицу поиска - структуру данных, 
которая позволяет дальнейшие преобразования проводить по уже известным соотношениям между пикселями. Это 
позволяет применять алгоритм для устранения искажений в реальном времени и реализовывать его на ПЛИС.
% TODO: точно можно?

\vspace{\baselineskip}


% Так как в используемой модели искажения хода луча считаются центрально симметричными и зависят только от его удаления от центра изображения, 
% рассмотрим ход падающего луча в координатах $(Z_c, \rho)$, изображённый на рисунке \ref{pic:scara_graph}.

% \addimghere{scara_graph}{0.5}{Нахождение обратной проекции для используемой модели}{pic:scara_graph} 

% Для модели (\ref{eqn:scaramuzza}) обратная проекция записывается как 
% \vskip 12pt
% \begin{equation}
%     \label{eq:back_scara}
%     \left[\begin{matrix}u_i\\v_i\\\end{matrix}\right] = \left[\begin{matrix} \frac{x_c}{\lambda}  \\  \frac{y_c}{\lambda} \\\end{matrix}\right],
% \end{equation} 
% \vskip 24pt
% где $\lambda = \rho_c / \rho_i$ - масштабный коэффициент. 

% Для нахождения $\rho_i$ был применён метод последовательных приближений. Блок-схема алгоритма, реализующего обратную проекцию, изображена  
% на рисунке  \ref{pic:newton_scheme}.  Сравнение изображений инструмента  калибровки и описанного алгоритма для центральной области изображения
% представлено на рисунке \ref{pic:central_pics}. 

% \addimghere{remapped_images}{0.8}{Изображения, скорректированные алгоритмом (слева) и MATLAB (справо)}{pic:central_pics} 


% \addimghere{flowchart}{0.5}{Блок-схема алгоритма обратной проекции}{pic:newton_scheme}         % Как происходит выпрямление
\subsection{Выводы по второму разделу}

Описан принцип устранения искажений сверхширокоугольных линз с выбором области интереса. Разработан алгоритм нахождения обратной 
проекции для fisheye-изображения. Описано устройство системы стереозрения. Разработана её виртуальная модель в среде Unity вместе с 
алгоритмом передачи изображений с виртуальной камеры в программу обработки изображений.  
Виртуальная модель системы позволяет перейти к её испытаниям.       % Выводы по разделу

\newpage
\section{Моделирование системы}
\subsection{Обоснование выбора ПО}
\label{sec:software}
Разработку и первоначальные испытания алгоритма стереозрения целесообразно проводить в виртуальной среде. Это позволяет значительно упростить разработку, 
так как уменьшает время на проверку гипотез и расходы на реальное оборудование, особенно в случае неудачных испытаний. Из-за этих факторов виртуальное моделирование 
в робототехнике приобрело широкое распространение и активно применяется, например, для разработки систем локализации и навигации беспилотного транспорта \cite{simulations}. 
Возросшее качество компьютерной графики к тому же позволило моделировать реалистичное окружение, что особенно важно при работе с системами технического зрения. 

Требованиями к виртуальной среде является возможность симулировать несколько широкоугольных камер и настраивать их параметры, 
легко интегрировать алгоритмы технического зрения и создать окружение, приближенное к тому, в котором может работать алгоритм в реальности. На данный момент исследователю доступен 
широкий выбор программного обеспечения, подходящего для этой задачи. В приложении А представлено сравнение имеющихся предложений по основным изложенным выше 
требованиям.

По результатам оценки собранные сведений принято решение проводить разработку в симуляторе Unity. Он позволяет подробно настраивать камеру и эмулировать сверхширокоугольные объективы 
с разными параметрами, строить реалистичные сцены благодаря свободному импорту моделей, а при программировании можно использовать сторонние программы в виде динамически подключаемых библиотек. 
По функционалу так же подходит NVIDIA Isaac Sim, но от него пришлось отказаться из-за высоких системных требований и новизны продукта.    

Разрабатываемое решение предполагается использовать в ПО робота, % TODO: не совсем ПО всё-таки
поэтому должно реализовываться на одном из популярных и быстродейственных языков программирования. Учитывая необходимость интеграции с Unity и потребность использовать популярные 
библиотеки, выбран язык C++.  Другим важным фактором является библиотека обработки изображений. В качестве основы для программной части была выбрана библиотека OpenCV, 
являющаяся стандартом при разработке систем технического зрения.        % Выбор ПО  
\subsection{Виртуальное моделирование системы}
\label{sysmodel}

Описанная система была смоделирована в среде Unity, её внешний вид представлен на рисунке \ref{pic:unity_model}. Мир 
Unity предназначен для базовой проверки работоспособности испытываемого принципа, поэтому не содержит подробной модели
какого-либо робота. В нём присутствуют: плоскость земли, кронштейн, на котором сверхширокоугольные камеры закреплены под 
углом $90^\circ$, подвижный калибровочный узор и объекты-цели, предназначенные для оценки расстояния. В силу особенностей 
работы многих алгоритмов стереозрения эти объекты должны сильнотекстурированы \cite{disparity_review}. 

\addimghere{unity_view}{0.7}{Внешний вид сцены в Unity}{pic:unity_model} 

В Unity создание сцены происходит с использованием встроенных примитивов, импортированных файлов моделей популярных форматов
или моделей из магазина ассетов, предлагающего обширную библиотеку объектов и текстур, повторяющих различные реальные объекты. В данной 
работе использовались как и примитивы для создания простых объектов, так и модели из магазина для имитации препятствий и окружения.

Для моделирования камеры "рыбий глаз" использовался аддон Dome Tools из магазина Unity Asset Store.    
Он позволяет моделировать сверхширокоугольные объективы с разным углом зрения в эквидистантной проекции \cite{dome_tools}. 
При этом с точки зрения других искажений, не относящихся к моделированию правильной проекции, снимки с этой камеры получаются идеальными. % FIXME: моделированию чего? 
Окно настроек камеры представлено на рисунке \ref{pic:camera_settings}.

\addimghere{camera_settings}{0.5}{Окно настроек виртуальной камеры в Unity}{pic:camera_settings} 

Для виртуальной камеры доступны настройки угла зрения и виньетки по краям изображения,  % FIXME: виньетки -> абберации ??
а также различных параметров рендеринга, влияющих на уровень  детализации получаемого изображения.      % Моделирование системы 
% \subsection{Передача изображений с виртуальных камер для обработки}

Unity в качестве основного языка программирования использует C\# и не позволяет импортировать 
выбранную библиотеку технического зрения OpenCV напрямую. Однако как уже упоминалось в секции 
\ref{sec:software}, данная среда позволяет интегрировать плагины, в том числе написанные на 
других языках программирования, в виде динамически подключаемых библиотек (DLL). Именно этот 
подход был использован для передачи изображений - была разработана библиотека с функциями для 
обработки снимков, которая компилировалась отдельно и затем импортировалась в Unity. 

Для реализации описанного в этой работе принципа и упрощения разработки в библиотеке были реализованы 
следующие функции:
\begin{itemize}
    \item initialize - выделяет память под нужные структуры, заполняет таблицу поиска и создаёт окна 
    для отображения будущих изображений.
    \item getImages - передаёт изображения в библиотеку и производит обработку изображений в 
    соответствии с аргументами. 
    \item takeScreenshot - производит ту же обработку входного изображения, что и предыдущая функция,
     но результат сохраняет в файл. 
    \item processImage - передаёт обратно в скрипт уже обработанные изображения для удобного отображения
    в интерфейсе Unity. 
    \item terminate - высвобождает память по завершении работы симуляции.      
\end{itemize}

Обработка, упомянутая в пунктах с функциями getImages и takeScreenshot, заключается в конвертации 
цветового пространства изображений (из RGBA, используемого в Unity, в BGR, используемый в OpenCV), 
горизонтальном зеркальном отображении (из-за разного положения точки отсчёта координат пикселей)
и, наконец, устранениb искажений в области интереса. Так как для наилучшей работы стереосопоставления
снимки должны быть синхронизированы (то есть получены камерами в один момент времени), обработка        % FIXME: ну не совсем одновременно и мб лучше это описать с другого ключа - типа функция может обрабатывать n снимков?
производится сразу для двух изображений. 

Также написан специальный скрипт, который управляет процессом взаимодействия симуляции и библиотеки. 
Эта программа имеет окно настроек, которое позволяет задать
используемые камеры, элементы управления параметрами  и выбрать область интереса с помощью задания углов.
 Внешний вид этого окна представлен на рисунке \ref{pic:connec_inter}.

\addimghere{connector_settings}{0.7}{Внешний вид окна настроек скрипта}{pic:connec_inter}

Сразу после запуска симуляции происходит запуск дополнительных вычислительных потоков для обработки 
кадров с каждой пары камер и передача настроек скриптом в библиотеку для построения таблиц поиска.           % FIXME: ну написано очень криво. Использование
Использование нескольких потоков позволяет не блокировать работу симуляции на время обработки изображений.  
Далее каждое обновление кадра скрипт считывает изображения с камер и помещает их в память соответствующего 
потока. Поток же работает независимо и обрабатывает каждую следующую пару изображений после готовности 
предыдущей.            % Про виртуальную камеру 
\subsection{Модели сверхширокоугольной камеры}
\label{camera_model}
Сложности, возникающие при использовании существующих алгоритмов стереозрения  в применении к сверхширокоугольным камерам, связаны с
 особенностями их оптической системы. Объективы  этих камер имеют в своей основе сложную систему линз, пример схемы которой вместе с примером
 получаемого изображения представлены на рисунке \ref{pic:fyscheme}. Особенности этой системы позволяют достигать очень высоких углов обзора,
  но также являются причиной аберрации и характерных искажений изображения. Чтобы описать свойства проецирования широкого набора таких
камер исследователи прибегают к аппроксимациям, называемым моделями камер. 

\addtwoimghere{fisheye_scheme}{0.45}{fisheye_example}{0.45}{Схема хода лучей объектива "рыбий глаз" (слева), пример изображения (справа)\cite{fy_exmp}}{pic:fyscheme}

Модель проекции для камеры это функция, которая описывает преобразование из точки трёхмерного пространства  в области зрения 
камеры ($P=[x_c, y_c, z_c]^T$) в точку на плоскости изображения ($p=[u, \nu]^T$), как показано на \ref{pic:fy_geom}. Единичная            % не совсем единичная 
полусфера $S$ с центром в точке $O_c$ на данной схеме описывает поле зрения. На ней также лежит точка $P_C$, являющаяся результатом обратной проекции.    %$\pi^{-1}_c({p})$
Угол $\theta$ является углом падения для рассматриваемой точки, а угол $\phi$ откладывается между положительным направлением оси $x$ и $O_{i}{p}$. 

Помимо самой проекции модели камер включают в себя описания нескольких типов искажений, накладываемых линзой. В сверхширокоугольных объективах самыми существенными являются 
радиальные - искажения, проявляющиеся сильнее по мере удаления от проекционного центра. Поэтому далее в этой секции модели будут рассматриваться именно с точки зрения 
описания радиальных искажений. 

\addimghere{projection_geometry}{0.5}{Схема проекции точки трёхмерного пространства в точку на изображении}{pic:fy_geom}

Перспективная проекция, которая обычно используется в качестве модели ортоскопической камеры, не способна спроецировать всё широкоугольное пространство на кадр 
конечного размера. Поэтому при описании и разработке fisheye-объективов опираются на другие виды проекций  \cite{projections}:

\begin{eqseries}
    \begin{equation}
        \label{fy1}
    r = 2 f tan(\theta/2),  
    \end{equation}
    \begin{equation}
        \label{fy2}
    r = f \theta,
    \end{equation}
    \begin{equation}
        \label{fy3}
    r = 2 f sin(\theta/2),
    \end{equation}
    \begin{equation}
        \label{fy4}
    r = f sin(\theta).
    \end{equation}
\end{eqseries}    

Но реальные искажения не всегда в точности следуют заданным уравнениями.
По этой причине fisheye-проекции выгоднее аппроксимировать другими функциями \cite{opencv_model}.

В настоящий момент есть несколько распространённых моделей, аппроксимирующих реальные искажения подобных объективов. Модель Канналы и 
Брандта \cite{opencv_model} для линз с радиально симметричными искажениями реализована в OpenCV и выражает их через угол падения луча света на линзу, а не расстояние  \pdfmargincomment{https://stackoverflow.com/questions/31089265/what-are-the-main-references-to-the-fish-eye-camera-model-in-opencv3-0-0dev}
от центра изображения до места падения, как это делалось в более ранних моделях. Авторы посчитали, что для описания типичных искажений достаточно 
пяти членов полинома с нечётными степенями. Таким образом, указанную модель можно записать следующими уравнениями:
\begin{eqseries}
    \begin{equation}	
        \delta r = k_1\theta + k_2\theta^3 + k_3\theta^5 + k_4\theta^7 + k_5\theta^9,
        \label{eqn:kannala_r}
    \end{equation}
    \begin{equation}	
        \begin{pmatrix}u\\v\end{pmatrix} = \delta r(\theta)\begin{pmatrix}cos(\phi)\\sin(\phi)\end{pmatrix},
        %\delta r = k_1\theta + k_2\theta^3 + k_3\theta^5 + k_4\theta^7 + ... + k_n\theta^{n+1}
        \label{eqn:kannala_uv}
    \end{equation}
\end{eqseries}
% TODO: indents
где $\theta$ - угол падения луча, определяемый выбранным типом проекции,

\qquad $\phi$ - угол между горизонтом и проекцией падающего луча на плоскость изображения, 

\qquad $r = sqrt(x^2+y^2)$ - расстояние от спроектированной точки до центра изображения, 

\qquad $f$ - фокусное расстояние, 

\qquad $k_1 \dots k_5$ - параметры модели.

Модель Мея \cite{mei} является более общей версией модели Гейера \cite{geyer} и позволяет использовать разные 
функции искажения для моделирования зеркал разоичного вида. Изначально она была создана для более 
эффективного моделирования катадиоптрических камер, но оказалась также весьма пригодной и для fisheye-камер. 

\begin{equation}
    \vspace{12pt}
    \begin{pmatrix}u\\v\end{pmatrix}=\left[\begin{array}{l}
	f_{x} \frac{x}{\alpha d+(1-\alpha) z} \\
	f_{y} \frac{y}{\alpha d+(1-\alpha) z}
	\end{array}\right]+\left[\begin{array}{l}
	c_{x} \\
	c_{y}
	\end{array}\right]
    \vspace{12pt}
\end{equation}

где $\alpha$ - параметр модели. 

Также большое распространение получила модель Скарамуззы \cite{scaramuzza}, которая легла в основу Matlab Omnidirectional 
Camera Calibration Toolbox. Она связывает точки на изображении с соответствующей им точкой в координатах камеры
следующим образом:
\vskip 12pt
\begin{equation}	
    \begin{pmatrix}X_c\\Y_c\\Z_c\end{pmatrix} = \lambda \begin{pmatrix}u\\v\\a_0 + a_2 r^2 + a_3 r^3 + a_4 r^4\end{pmatrix},
    %\delta r = k_1\theta + k_2\theta^3 + k_3\theta^5 + k_4\theta^7 + ... + k_n\theta^{n+1}
    \label{eqn:scaramuzza} 
\end{equation}
\vskip 24pt

где $a_0 ... a_4$ - коэффициенты, описывающие параметры модели,

\qquad $\lambda$ - масштабный коэффициент.

Существуют и менее распространённые модели, не использующие полиномы для описания искажений. Одной из них является
модель двух сфер \cite{double_sphere}. 
Она находит положение пикселя, проектируя точку в несколько этапов - сначала на первичную сферу, потом на вторую 
сферу меньшего диаметра и смещённую на расстояние $\xi$. Наконец точка проецируется на  плоскость изображения,
 сдвинутую на расстояние $\frac{\alpha}{1-\alpha}$ относительно центра второй сферы. Модель проекции представлена 
на рисунке \ref{pic:ds_model}. Таким образом, для описания радиальных искажений достаточно всего двух параметров.
 Модель также реализована в нескольких популярные программ для калибровки камер (Basalt, Kalibr). 
 Записывается она следующим образом:
\begin{eqseries}
    \begin{equation}	
        d_1 = \sqrt{x^2+y^2+z^2}, 
    \end{equation}
    \begin{equation}	
        d_2 = \sqrt{x^2+y^2+(\xi*d_1+z)^2 }, 
    \end{equation}
    \begin{equation}	
        \begin{pmatrix}u\\v\end{pmatrix} = \begin{pmatrix}f_x * \frac{x}{\alpha*d_2+(1-\alpha)(\xi*d_1+z)} \\
                                                    f_y * \frac{y}{\alpha*d_2+(1-\alpha)(\xi*d_1+z)} \end{pmatrix}.
        %\delta r = k_1\theta + k_2\theta^3 + k_3\theta^5 + k_4\theta^7 + ... + k_n\theta^{n+1}
        \label{eqn:ds}
    \end{equation}
\end{eqseries}
\addimghere{double_sphere}{0.7}{Модель  двух сфер}{pic:ds_model}

              % Обзор моделей камер

В этом разделе рассматриваются модели линз в их применении к описанной системе стереозрения. 
Так как классические методы стереозрения, требующие ректификации, всё ещё остаются самыми доступными и производительными \cite{disparity_review}, 
необходимо выбрать модель искажений, способную обеспечить систему наиболее точно восстановленными изображениями. 
Результатом работы алгоритма стереозрения является карта глубины воспринимаемого пространства. Подобные карты могут
использоваться, например, алгоритмами навигации и локализации для построения карты окружения робота. 
Исходные данные сравниваемых моделей представлены в приложении \ref{app-b}.

Точность работы 
стереокамер зависит от того, насколько точно карта глубины передаёт реальную информацию о форме объектов. Оценить 
эту характеристику для отдельной стереопары проблематично, так как она зависит от множества факторов. Поэтому оценка 
качества работы системы с различными моделями проведена в сравнении с  эталонами.

Кроме указанных моделей в сравнениях можно увидеть ATAN. Эта модель представляет из себя идеальную эквидистантную fisheye-проекцию, 
которая заложена в использованную виртуальную широкоугольную камеру. 
Она добавлена к сравнению как эталонный способ устранения искажений из предположения, что с ней результаты стереосопоставления будут 
самыми лучшими. Использование её на реальных объективах ограничено отсутствием возможностей по калибровке. 

Кроме того, виртуальная сцена позволяет разместить сразу несколько объектов в одной точке пространства, таким образом возможно
в существующую модель в Unity добавить ещё 2 камеры, совпадающие по параметрам и положению с виртуальными камерами на 
рисунке \ref{pic:2cam_scheme}. Разрешение новых камер выбрано исходя из размеров проекции $\nu_i$ области интереса на широкоугольном
снимке. Для камер "рыбий глаз" с разрешением $1080*1080$ пикселей она составляет 287482 пикселей, что аналогично камере 
с разрешением $540*540$. Камеры полностью подчиняются перспективной проекции (\ref{eq:uv_to_xyz}) и не содержат каких либо оптических искажений.
 Далее эти камеры будут называться эталонными камерами, а стереопара, которую они составляют - эталонной стереопарой (обозначена на графиках как REG).

Наиболее сильные искажения на изображении "рыбий глаз" возникают по краям изображения, поэтому исследуемая виртуальная стереопара построена так,
чтобы зона пересечения полей зрения камер располагалась на максимальном удалении от центра. Это соответствует ортогональному расположению
камер с углом между принципиальными осями в $90^\circ$.  При этом объективы камер имеют одинаковые параметры, поле зрения в $180^\circ$,
 а расстояние между ними составляет 20см. %Описанная конфигурация изображена на рисунке \ref{pic:2cam_scheme}.

% Здесь область пересечения полей зрения камер $C_0$ и $C_1$ обозначенная красным.
% Эта область эквивалентна области пересечения полей зрения двух камер с полями зрения $90^\circ$ (обозначены
% оранжевым), повёрнутых на $\beta_0 = \beta_1  = 45^\circ$ в сторону области интереса. Тогда $B$ - база стереопары.

\subsection{Экспериментальное исследование моделей}
\label{model_exps}

С помощью этой стереопары были получены снимки калибровочного узора шахматной доски, 
которые потом  направлены в соответствующие программные пакеты для нахождения параметров каждой модели. 
Параметры моделей занесены в разработанную библиотеку для устранения искажений на снимках виртуальной 
плоскости на расстояниях от 1 до 10м с шагом 1м. 
Пример снимка после устранения искажений и снимка с соответствующей эталонной камеры показан на рисунке \ref{pic:2images_compar}.

\addtwoimghere{pic_fy}{0.4}{pic_reg}{0.4}{Слева - снимок после устранения искажений (модель Скарамуззы); справа - эталонное изображение}{pic:2images_compar}

Визуально снимки очень похожи, однако при более детальном рассмотрении на левом изображении заметна меньшая чёткость,
вероятно, связанная с интерполяцией пикселей в процессе проецирования. Также присутствуют малозаметные искажения геометрии. Более 
явно эти дефекты для разных моделей можно увидеть на разностном изображении, представленном на рисунке \ref{pic:difference}. 

\addimghere{diffs}{0.9}{Разностные изображения (резкость увеличена). Чем светлее участок, тем сильнее различия}{pic:difference}

Визуальный анализ этих изображений подтверждает наблюдения и показывает малозначительность отличий для большинства моделей. 
Это позволяет применять изображения, полученные описанным путём, в стереосопоставлении. 

Оценка качества оценки глубины произведена на основе измерения среднеквадратичного отклонения и дисперсии полученных точек 
глубины от их предполагаемой позиции. Виртуальность эксперимента позволяет точно знать положение исследуемого объекта и, соответственно, 
точно определять ошибку. Как уже упоминалось, эффективность стереосопоставления \cite{SGBM} зависит от текстуры наблюдаемого объекта, 
поэтому для устранения влияния этого фактора эксперименты были проведены с различными текстурами. Использованный алгоритм стереосопоставления
опирается на поиск отклика на участок одного изображения на эпиполярной линии другого. Таким образом, текстура с одинаковыми повторяющимися 
фрагментами может привести к ложным срабатываниям и карте глубины низкого качества. Поэтому в экспериментах использовались три текстуры: часто 
повторяющийся узор обоев; текстура, состоящая как из повторяющихся с разной частотой, так и из случайных элементов; и полностью неповторяющееся изображение.
Использованные текстуры приведены на рисунке \ref{pic:textures}. Графики в этом и следующем разделе отображают усреднённые результаты по всем текстурам, отдельные графики приведены в приложениях. 

\addimghere{pics/textures}{0.7}{Пример использованных текстур. а)~переодическая; б)~смешанная; в)~непереодическая}{pic:textures}

По полученным картам глубины можно осуществить 3D-реконструкция сцены. Результата реконструкции представлен  в виде облака точек, 
 аналогичного изображённому на рисунке \ref{pic:raw_pointcloud}. По соответствующим этому снимку дистанциям строится модель целевой плоскости, 
 а точки за пределами её окрестности отбрасываются. Далее из оставшихся выбираются 1000 случайных точек, что позволяет считать ошибку с 
 одинаковой точностью для всех расстояний.  % FIXME: размеры окрестности и точно ли _точность_ одинаковая? 
Эти точки используются для вычисления ошибки. 

\addimghere{pointcloud}{0.6}{Неочищенное облако точек}{pic:raw_pointcloud}

Результаты оценки качества нахождения глубины приведены на рисунке \ref{pic:quality}. По оси абсцисс отложено расстояние до исследуемого объекта,
по оси ординат среднеквадратичное отклонение. На графике так же отмечена линейная аппроксимация самой точной модели.
 На рисунке \ref{pic:mean} изображена зависимость матожидания дистанции до поверхности от расстояния
до исследуемой плоскости. Вертикальные линии показывают разброс точек.  Каждая ломанная соответствует своей модели и усреднена по всем текстурам.

\addimg{pics/depth_quality}{1}{Точность построения карты глубины}{pic:quality}

\addimg{pics/distance}{1}{Оценка дистанции до плоскости}{pic:mean}

Как можно заметить по графикам, наилучший результат в оценке глубины логичным образом показывает эталонная стереопара. Далее следует 
идеальная модель, заложенная в виртуальную камеру, демонстрируя, что даже в случае полного соответствие прямой и обратной проекций часть 
информации в изображении теряется или искажается. Из исследуемых моделей самую низкую ошибку демонстрируют модель Мея с результатом
0.04м на метр удаления и модель Канналы-Брандта с результатом 0.05м/м, модель Скарамуззы показывает себя хуже других с ошибкой в 0.06м на метр удаления.
 Кроме того, по второму графику заметно, что менее точные моделей имеют тенденцию к недооценке расстояния. 

Так же можно рассмотрено распределение значений глубины точек, попавших в конечную выборку. При более точных измерениях
распределение близко к нормальному, в то время как ухудшение точности в силу повышения расстояния или использования другой модели
сопровождается возникновением дополнительных пиков или скосов на гистограмме, как показано на рисунке \ref{pic:histo_virt}. 

\addimghere{pics/distributions}{1}{Распределение точек глубины. а)~Модель Мея на 1м; б)~Модель Канналы-Брандта на 9м; в)~Модель Скарамуззы на 5м}{pic:histo_virt}




\subsection{Выводы по второму разделу}

Проведённые исследования показали, что в применении к стереозрению наилучшие результаты демонстрирует модель Канналы и Брандта. Она демонстрирует наименьший 
прирост ошибки, который составляет в среднем 0.06м на метр удаления от камеры, что на 27\% больше, чем для эталонной стереопары. Кроме того, модель является довольно доступной 
в плане автоматической калибровки и доступной в нескольких библиотеках технического зрения. Несмотря на меньшую точность, 
стереосистема, основанная на камерах со сверхширокоугольными объективами, может успешно применяться \cite{}.

В этот раздел планируется добавление ещё одного эксперимента по оценке степени "скашивания" облаков точек. 

% \newpage
% \section{Исследование точности калибровочных моделей сверхширокоугольных объективов}

\newpage
\section{Экспериментальное исследование системы стереозрения}

Предлагаемая стереосистема показала свою работоспособность и удовлетворительные результаты в виртуальных тестах 
точности. Это позволяет перейти к испытаниям системы на реальных камерах. 

Для этого использованы 2 имеющиеся в наличии камеры % TODO: написать тип + ссылки
с линзами 1.45mm F2.2 1/1.8 FOV $190^\circ$ (AC123B0145IRM12MM), закреплённые под углом $90^\circ$ в корпусе,
полученным 3D-печатью. База стереопары составляет в таком случае примерно  $72мм$. Изображение модуля системы 
стереозрения представлен на рисунке \ref{pic:cam_case}. 

\addimghere{pics/case_photo}{1}{Система стереозрения в корпусе}{pic:cam_case}

Пример, пары fisheye-изображений, выдаваемых таким модулем камер представлен на рисунке \ref{pic:stere_fy_img}. 
Как можно заметить, у этих камер круг линзы не полностью вписан в кадр, а её центр смещён относительно центра кадра,
 что уменьшает количество полезной площади кадра, доступной для устранения искажений. Кроме того, заметна виньетка 
по краям, которая может мешать поиску соответствий. Всё это повышает вклад сенсоров в общую ошибку. 

\addimghere{pics/real_fy_shots}{1}{Снимки с двух реальных камер с объективами "рыбий глаз"}{pic:stereo_fy_img}

Ошибка условий измерений минимизируется использованием искусственного освещения, расстояние до целевого объекта
контролируется разметкой на экспериментальном столе, нанесённой с помощью рулетки. Аналогично виртуальным испытаниям 
использованы 3 текстуры, распечатанные на листе матовой бумаги, для уменьшения влияния свойств наблюдаемой поверхности.

Каждая камера по-отдельности откалибрована по снимкам узора шахматной доски, полученные параметры параметры моделей 
представлены в таблице \ref{}. Примеры изображений с устранёнными искажениями для каждой модели представлены на рисунке
 \ref{pic:dewarped_exmples_real}. Заметно что на них присутствуют пустоты, вызванные усечённым форматом оригинального 
 изображения. Тем не менее, эти пустоты не мешают дальнейшему эксперименту, так как оставляют достаточное поле зрения 
 незатронутым.

\addimg{pics/dewarped_real}{1}{Примеры изображений с устранёнными искажениями}{pic:dewarped_exmples_real} 
\vspace{\baselineskip}
\subsection{Экспериментальное исследование реальной системы стереозрения}

Экспериментальная установка состоит из модуля системы стереозрения, закреплённого на краю стола, 
разметки с шагом 25см и целевого объекта в виде коробки с приклеенным изображением текстуры. Снимок 
экспериментальной установки представлен на рисунке \ref{pic:table_photo}. 

С помощью этой установки получены снимки целевого объекта с тремя разными текстурами, вручную помещённого на
удаление от 0.25м до 1.5м с шагом 0.25м. Они переданы в тот же алгоритм подсчёта ошибки, описанный в \ref{}. 

Результаты оценки качества нахождения глубины приведены на рисунке \ref{pic:quality_real}. По оси абсцисс 
отложено расстояние до исследуемого объекта, по оси ординат среднеквадратичное отклонение. На рисунке \ref{pic:mean_real}
 изображено зависимость матожидания дистанции до поверхности от расстояния до исследуемой плоскости и дисперсия точек.
Каждая ломанная соответствует своей модели и усреднена по всем текстурам.

\addimghere{pics/depth_quality_real}{1}{Точность построения карты глубины}{pic:quality_real}

\addimghere{pics/distance_real}{1}{Оценка дистанции до плоскости}{pic:mean_real}



% Рассмотрим работу модуля кругового стереозрения на основе четырёх камер, аналогичного 
% изображённому на рисунке \ref{pic:ssvs}. Применение описанного ранее алгоритма позволяет
% использовать каждую камер сразу в двух стереопарах, обеспечивая широкое покрытие.
% Подобный модуль мог бы обеспечить робота как полным панорманым обзором, так и почти полным панорамной оценкой глубины. 

% Сборка виртуальной модели системы произведена в Unity. Модуль помещён в одну из локаций 
% и использован для получения снимков со всех камер. По полученным картам глубины построено
% общее облако точек, изображённое вместе с исходной локацией на рисунке \ref{pic:module_cloud}. 

% \addimghere{unity_with_ptc}{1}{Слева - локация, в которыой были сделаны снимки; справа - результат реконструкции (обрезан по высоте)}{pic:module_cloud}
% % \addtwoimghere{unity_topdown}{0.4}{white_ptc}{0.4}{Слева - локация, в которыой были сделаны снимки; справа - результат реконструкции (обрезан по высоте)}{pic:module_cloud}

% В облаке заметны очертания окружающих поверхностей. Дальнейшие испытания целесообразно проводить 
% с использованием одного из распространённых пакетов автономной локализации и построения карты в рамках.
\vspace{\baselineskip}
\subsection{Заключение по пятому разделу}

Система стереозрения готова к испытаниям в системах навигации. Необходимы виртуальные 
тесты с применением, например, ROS.

Этот раздел будет дополняться результатами экспериментов по навигации. Кроме того, проходят 
подготовку реальные камеры для оценки точности построения карты глубины по методике \ref{model_exps},
но уже в настоящих условиях.

% Описанная в главе \ref{sysmodel} виртуальная модель системы стереозрения позволяет проводить 
с ней испытания для оценки работоспособности и сравнения с аналогами в контролируемой среде.           

Результатом работы алгоритма стереозрения является карта глубины воспринимаемого пространства. Подобные карты могут
использоваться, например, алгоритмами навигации и локализации для построения карты окружения робота. Точность работы 
этих алгоритмов зависит от того, насколько точно карта глубины передаёт реальную информацию о форме объектов. Оценить 
эту характеристику для отдельной системы стереозрения проблематично, так как она зависит от множества факторов.         % TODO: дописать факторы или сослаться на статью
Поэтому оценка качества работы системы проведена в сравнении с виртуальной моделью традиционной стереопары.

Здесь в виртуальных экспериментах для сравнения так же применяется эталонная стереопара.  % TODO: чувствуется какая-то недосказанность
       % План экспериментов
% 
\subsection{Оценка качества устранения искажений}

Добавление в виртуальную модель эталонной камеры позволяет оценить качество устранения искажений путём сравнения
её снимков со снимками виртуальной камеры-обскуры модуля устранения искажений. % FIXME: никакой "модуль..." ранее не вводился, нужно переписать

Параметры полинома (\ref{eqn:scaramuzza}) для устранения искажений получены с помощью MATLAB Camera Calibrator и занесены 
в код библиотеки обработки изображений. % FIXME: опять-таки будто про OpenCV речь, надо как-то ссылаться на свой алгоритм\модуль
Снимки после устранения искажений и с эталонной камеры показаны на рисунке \ref{pic:2images_compar}.

\addtwoimghere{pic_fy}{0.4}{pic_reg}{0.4}{Слева - снимок после устранения искажений; справа - эталонное изображение}{pic:2images_compar}

Визуально снимки очень похожи, однако при более детальном рассмотрении на левом изображении заметна большая резкость, вызванная, 
вероятно, округлениями чисел в процессе проекции. Также присутствуют малозаметные искажения геометрии. Более 
явно эти дефекты можно увидеть на разностном изображении, представленном на рисунке \ref{pic:difference}. 

\addimghere{difference}{0.5}{Разностное изображение (резкость увеличена). Чем светлее участок, тем сильнее различия}{pic:difference}

Визуальный анализ этого изображения подтверждает различия в резкости обеих картинок и наличие искажений преимущественно в левой части снимка. Тем не менее, искажения
достаточно несущественны и проявляются лишь близь краёв исходного изображения, что позволяет считать подобные снимки пригодными 
к применению в системе стереозрения.           % Сравнение картинок
% 
\subsection{Исследование точности оценки глубины}  % понятие облака точек не введено, поверхность какая? 

В качестве целевой поверхности выбрана виртуальная плоскость с нанесённой на неё нерегулярной текстурой высоко разрешения. Размеры и
положение плоскости относительно камер известно с высокой точностью, что позволяет сравнить результаты стереореконструкции
с реальным положением целевого объекта и оценить ошибку. В качестве метрики оценки выбрано среднее квадратичное отклонение
положения точек от модели плоскости. 

Исследования проведены в MATLAB, снимки получены с помощью виртуальной модели при  одинаковых условиях освещённости и 
постоянных  настройках всех  компонентов системы.  Ширина базы $B$ выбрана равной 1  метру.  Сначала для обеих стереопар 
проведена калибровка по снимкам узора шахматной доски с помощью методики \cite{stereo_calib}, что позволило получить 
внутренние и внешние параметры камер стереопары. % FIXME: ну вообще они у нас как бы есть. Мб сравнить ещё оценки с реальными 
Далее происходит построение карты расхождений методом полу-глобального  сопоставления \cite{SGBM} и 3D-реконструкция сцены. 
Результата реконструкции представлен  в виде облака точек,  изображённого на рисунке \ref{pic:raw_pointcloud}.

\addimghere{pointcloud}{0.7}{Неочищенное облако точек}{pic:raw_pointcloud}
\pdfmargincomment{Три вида с одного раркурса}
% TODO: сюда можно и картинку после обработки поинтклауда \addtwoimghere{}{}{}{}{}
По заданным параметрам строится модель целевой плоскости, а точки за пределами её окрестности отбрасываются. % FIXME: размеры окрестности
Затем скрипт перебирает все оставшиеся точки и рассчитывает среднее квадратичное отклонение по длине нормали от 
точки к плоскости. 
Описанный алгоритм применён к снимкам целевой плоскости на разных расстояниях с эталонной и исследуемой стереопар.
Результаты приведены на графике \ref{plot:MSE_compar}. На графике приведены значения среднеквадратичной ошибки и
дисперсия для исследуемой и эталонной стереопар при разном удалении целевого объекта.

\addimghere{alt1mErrorGraph}{1.0}{График среднеквадратичной ошибки и  дисперсии в зависимости от расстояния для эталонной и исследуемой стереопар}{plot:MSE_compar}

Как видно из аппроксимирующих прямых на графике, среднеквадратичная ошибка оценки поверхности в исследуемой системе
 в 3 раза больше этого показателя для традиционной стереопары. Меньшая точность исследуемой системы связана с влиянием 
на алгоритм стереосопоставления дефектов изображения, описанных в предыдущей секции. Кроме того, с увеличением расстояния
уменьшается разрешение, приходящееся на поверхность, и, соответственно, количество точек с подсчитанной глубиной. 
Тем не менее она остаётся в пределах $3\%$ от реального расстояния.  % FIXME: хммммм, как-то уж больно хорошо. Или  плохо... 10 см на 5 метрах


%         % Оценка равномерности поинтклауда стены
% \subsection{Выводы по третьему разделу}

Исследованы изображения, полученные с помощью виртуальных камер поле устранения искажений. Выполнено их сравнение 
с эталонными изображениями,  которое продемонстрировало пригодность снимков для использования в системах стереозрения.

В виртуальную среду добавлена эталонная стереопара для сравнения с разработанной системой. Проведён эксперимент,
по оценке качества облака точек, полученного с помощью предлагаемого решения, в сравнении  традиционной системой.
Полученные результаты сравнимы с актуальными пассивными система стереозрения \cite{passive_perf}, что позволяет перейти
к дальнейшим этапам реализации системы стереозрения на основе ортогонально расположенных камер. . 


  % Выводы по разделу

\newpage
\anonsection{Заключение}

В ходе работы рассмотрены существующие системы стереозрения, в особенности использующие информацию
со сверхширокоугольных камер. Изучены модели, описывающие подобные камеры. Рассмотрено ПО для виртуального
 моделирования робототехнических комплексов и систем технического зрения. Выбрано ПО для моделировани и 
разработки системы стереозрения, использующей объективы "рыбий глаз". 

Разработан и описан алгоритм устранения искажений в области интереса и его математическая модель.   % пафосно, конечно.  Так ли хорошо она описана?
Описана возможная конфигурация системы стереозрения, рассмотрен её принцип работы. Система 
 смоделирована в среде Unity. В процессе размещены и настроены камеры, подготовлены объекты для калибровки и 
дальнейших испытаний системы. Реализована передача изображений с виртуальных камер для обработки алгоритмами
компьютерного зрения. С помощью снимков из виртуальной модели исследовано качество устранения искажений и выполнено
 сравнение точности оценки глубины предлагаемой системой и варианта с обычными камерами. Результаты позволяют 
 перейти к исследованию работы системы с применением реальных камер.  % FIXME: "указывают на доработки" такое себе 

Дальнейшая работа будет сконцентрирована на оптимизации и повышении точности алгоритма устранения искажений,
разработке способа автоматической установки параметров системы при разных конфигурациях расположения камер. 
Эти этапы позволят исследовать применимость системы стереозрения при  разных схемах размещения камер и разной 
длине базы, а также рассмотреть качество построения карты и локализации с применением описанной 
системы сначала с виртуальным, а затем и с реальном роботом.   %Заключение-заключение

\newpage
\anonsection{СПИСОК ИСПОЛЬЗОВАННЫХ ИСТОЧНИКОВ}
\renewcommand{\refname}{}
\patchcmd{\thebibliography}{\section*{\refname}}{}{}{}
\patchcmd{\thebibliography}{\addcontentsline{toc}{section}{\refname}}{}{}{}
\patchcmd{\thebibliography}{\begin{itemize}}{\begin{itemize}[leftmargin=30mm]}{}{}
    \makeatletter
    \renewcommand*{\@biblabel}[1]{\hfill#1}
    \makeatother
{\def\section*#1{}
\bibliography{refs.bib} 
}
\newpage
% Приложения
\appsection{Приложение А} \hypertarget{app-a}{\label{app-a}}

\centering{Сравнение ПО для симуляции}


    \begin{table}[h!]             % TODO: Сыровато. Надо дополнить. Проверить ГОСТовость подписи  FIXME: Таблица не умещается, но вращать не хочется. 
        
        %\caption{Сравнение ПО для симуляции }
        \label{tab:sims}
        \rotatebox{90} {\begin{tabular}{|l|l|l|l|l|}
        \hline
        \textbf{Название} & \textbf{\begin{tabular}[c]{@{}l@{}}Симуляция \\ fisheye-камер\end{tabular}} & \textbf{\begin{tabular}[c]{@{}l@{}}Реалистичное \\ моделирование\end{tabular}} & \textbf{Интеграция кода}                                            & \textbf{Доступность}  \\ \hline
        Gazebo            & Возможна                                                                    & Возможно                                                                     & \begin{tabular}[c]{@{}l@{}}Возможна \\ посредством ROS\end{tabular} & Бесплатно               \\ \hline
        RoboDK            & Нет                                                                         & Затруднено                                                                     & Нет                                                                 & От 145€                \\ \hline
        Webots            & Затруднена                                                                  & Возможно                                                                       & Возможна                                                            & Бесплатно              \\ \hline
        CoppeliaSim       & Затруднена                                                                  & Затруднено                                                                     & Возможна                                                            & Бесплатно              \\ \hline
        NVIDIA Isaac Sim  & Возможна                                                                    & Возможно                                                                       & Возможна                                                            & Бесплатно              \\ \hline
        CARLA             & Затруднена                                                                  & Возможно                                                                       & Возможна                                                            & Бесплатно              \\ \hline
        Unity             & Возможна                                                                    & Возможно                                                                       & Возможна                                                            & Бесплатно              \\ \hline
        \end{tabular}}

    \end{table}
\clearpage

\appsection{Приложение Б} \hypertarget{app-a}{\label{app-b}}

\centering{Исходные данные моделей]}

\begin{table}[h!]
    \centering
    % \caption{Исходные данные}
    \label{tab:models_params}
    \rotatebox{90} {\begin{tabular}{|l|l|l|l|}
        \hline
    Модель         & ПО калибровки & \begin{tabular}[c]{@{}l@{}}Ошибка\\ калибровки\end{tabular}   & Параметры модели \\ \hline
    Мей            & OdoCamCalib   & 0.102               & $ \xi=1.47431, p_1=0.000166, p_2=0.00008, k_1=-0.208916, k_2=0.153247 $                 \\ \hline
    Каннала-Брандт & OdoCamCalib   & 0.099               & $ k_2=0.000757676, k_3=-0.000325907, k_4=0.0000403, k_5=-0.000001866, f_x=343.57, f_y=343.37$     \\ \hline
    Скарамузза     & MATLAB        & 0.121               & $ a_0=345.1319, a_1=-0.0011, a_2=5.7623 * 10^-7, a_3=-1.3985 * 10^-9 $               \\ \hline
    \end{tabular} }
\end{table}
\clearpage

% \appsection{Приложение Б} \hypertarget{app-b}{\label{app-b}}

% \centering{Программный код файла FisheyeDewarper.cpp}
% \begin{english}
% \lstinputlisting[language=C++, numbers=left]{code/FisheyeDewarper.cpp}
% \end{english}

% \clearpage

% \appsection{Приложение В} \hypertarget{app-c}{\label{app-c}}

% \centering{Программный код файла unity\_plugin.cpp}
% \begin{english}
% \lstinputlisting[language=C++, numbers=left]{code/unity_plugin.cpp}
% \end{english}

% \clearpage

% \appsection{Приложение Г} \hypertarget{app-d}{\label{app-d}}

% \centering{Программный код файла Connector.cs}
% \begin{english}
% \lstinputlisting[language=C, numbers=left]{code/Connector.cs}
% \end{english}

% \clearpage

% \appsection{Приложение Д} \hypertarget{app:matlab}{\label{app:matlab}}

% \centering{Программный код файла DisparityPlayground.m}
% \begin{english}
% \lstinputlisting[language=Matlab, numbers=left]{code/DisparityPlayground.m}
% \end{english}

% \clearpage
 % Код скрипта

% \lof
% \lot

\end{document} 
%%% Конец документа
